\begin{proof}\leavevmode
\paragraph{1. One uniform torsion annihilator.}
By the non-CM open-image theorem \cite{serre}, choose $n\ge2$ such that
$1+2^nM_2(\mathbf Z_2)\subset\rho_T(G_{\mathbf Q})$, and choose
$g\in G_{\mathbf Q}$ with $\rho_T(g)=uI$, $u=1+2^n$.
For an actual complex conjugation $\tau$ set
\[
 z=g^2\tau g^2\tau^{-1},\qquad e=v_2(u^4-1)=n+2.       \tag{H1}
\]
The element $g^2$ fixes every quadratic field. Its image in the
abelian group $\operatorname{Gal}(H_d/K)$ is inverted by $\tau$,
for every $d$ and every allowed $K$. Thus $z$ fixes every such
$H_d$, while $\rho_T(z)=u^4I$. If $P\in E(H_d)[2^\infty]$,
then $(u^4-1)P=0$, so
\[
                       2^e E(H_d)[2^\infty]=0.          \tag{H2}
\]
This bounds the whole invariant torsion group, with no restriction
on the number of conductor primes. It also holds for a compositum
of the ring-class fields in a finite construction, since $z$ fixes
each of them.

For the matrix-span statement used elsewhere, square preimages of
$I+2^nE_{11}$, $I+2^nE_{12}$ and $I+2^nE_{21}$. The squares belong
to every quadratic subgroup, and their images, together with $I$,
are
\[
 I,\quad I+(2^{n+1}+2^{2n})E_{11},\quad
 I+2^{n+1}E_{12},\quad I+2^{n+1}E_{21}.
\]
Their $\mathbf Z_2$-span contains $2^{n+1}M_2(\mathbf Z_2)$:
subtract $I$, use that $1+2^{n-1}$ is odd, and obtain $E_{22}$
from $I-E_{11}$. This bound also depends only on $E$.
